\documentclass[11pt]{amsart}

\usepackage[margin=1in]{geometry}
\usepackage{amsmath,amssymb,amsthm,mathtools}
\usepackage{hyperref}
\usepackage[shortlabels]{enumitem}

\newtheorem{theorem}{Theorem}[section]
\newtheorem{proposition}[theorem]{Proposition}
\newtheorem{lemma}[theorem]{Lemma}
\newtheorem{corollary}[theorem]{Corollary}

\theoremstyle{remark}
\newtheorem{remark}[theorem]{Remark}

\newcommand{\C}{\mathbb{C}}
\newcommand{\N}{\mathbb{N}}
\newcommand{\abs}[1]{\lvert #1 \rvert}
\newcommand{\norm}[1]{\lVert #1 \rVert}
\DeclareMathOperator{\dist}{dist}
\newcommand{\dd}{\mathrm{d}}
\newcommand{\ee}{\mathrm{e}}
\newcommand{\ii}{\mathrm{i}}
\newcommand{\rhoFn}{\rho}
\newcommand{\Fd}{\mathcal{F}_d}
\newcommand{\Sphere}{S^{2d-1}}
\newcommand{\Ann}[1]{\mathcal{A}_{#1}}

\title{Local Stable Phase Retrieval in Bargmann--Fock Space on $\C^d$}
\author{}

\begin{document}

\begin{abstract}
We give a higher-dimensional analogue of the verified one-dimensional
argument recorded in the source tree
\texttt{Constant/Latex/FockSPR}.  The main analytic task is to prove the
orthogonal coercivity estimate in $\Fd(\C^d)$.  We first establish it for
holomorphic polynomials with vanishing constant term.  The key observation
is that the circle estimates from
\texttt{Constant/Latex/FockSPR/CircleEstimates.tex} lift directly to the
unit sphere $\Sphere$ by averaging the global phase action
$\omega \mapsto \ee^{\ii t}\omega$.  After that, the annulus decomposition
and leakage estimates from
\texttt{Constant/Latex/FockSPR/BlockDecomposition.tex} go through with the
total degree replacing the one-variable degree.  A density argument then
extends the orthogonal estimate from polynomials to all of $\Fd$.  Finally,
we record the orthogonal reduction explicitly in the present notation and
verify phase retrieval at the constant function in $\Fd(\C^d)$.  This yields
a quantitative small-data stability theorem at the constant, with a
constant depending only on~$d$.  We also isolate a separate
compactness-based local-to-global passage, but we treat that final step as
an optional soft consequence rather than as part of the quantitative core.
\end{abstract}

\maketitle

\section{Statement}

For a measurable function $f \colon \C^d \to \C$, define
\[
\norm{f}_{\Fd}^2
:=
\frac{1}{\pi^d}\int_{\C^d} \abs{f(z)}^2 \ee^{-\abs{z}^2}\,\dd m(z),
\]
and let $\Fd = \Fd(\C^d)$ be the closed subspace of entire functions with
finite norm.
Throughout, $\langle \cdot,\cdot\rangle_{\Fd}$ denotes the Fock-space inner
product
\[
\langle f,g\rangle_{\Fd}
:=
\frac{1}{\pi^d}\int_{\C^d} f(z)\overline{g(z)} \ee^{-\abs{z}^2}\,\dd m(z),
\]
which is linear in the first slot.
By abuse of notation, $\norm{\cdot}_{\Fd}$ also denotes this ambient
Gaussian weighted $L^2$ norm on measurable functions; in particular,
expressions such as $\norm{\rhoFn \circ U}_{\Fd}$ are interpreted in that
sense.

As in the one-dimensional note, write
\[
\rhoFn(w) := \bigl\lvert\, \abs{1+w} - 1 \,\bigr\rvert.
\]
The one-dimensional note writes $\rho_0(w) := \abs{1+w}-1$ without the
outer absolute value.  In the argument below we only use that
\[
\abs{\rho_0(w)} = \rhoFn(w)
\qquad\text{and}\qquad
\abs{\rhoFn(w)-\rhoFn(z)} \le \abs{w-z}.
\]
Consequently, every estimate from the one-dimensional note whose statement
depends only on the $L^2$ norm of $\rho_0 \circ P$ transfers unchanged to
$\rhoFn \circ P$, because for every measurable $P$ one has
\[
\norm{\rho_0 \circ P}_{L^2} = \norm{\rhoFn \circ P}_{L^2}.
\]
Separately, the same $1$-Lipschitz bounds remain available.  We do not use
any signed identity for $\rho_0$ after passing to $\rhoFn$.

The real analytic content is the following orthogonal estimate.

\begin{theorem}[Orthogonal coercivity on $\C^d$]\label{thm:orthog}
Fix $d \ge 1$.  There exists a constant $C_d > 0$ such that every
$U \in \Fd$ with $U(0) = 0$ satisfies
\[
\norm{U}_{\Fd}^2
\le
C_d^2\, \norm{\rhoFn \circ U}_{\Fd}^2.
\]
\end{theorem}

\noindent
Theorem~\ref{thm:orthog}, combined with the orthogonal reduction stated
explicitly in Proposition~\ref{prop:orth-reduction}, implies the following
small-data stability statement, which is the main quantitative consequence
of the note.

\begin{corollary}[Small-data stability at the constant]\label{cor:small}
Fix $d \ge 1$, and let $C_d$ be as in Theorem~\ref{thm:orthog}.  Then there
exist
\[
\delta_d := \frac{1}{C_d + 1},
\qquad
\widetilde M_d := 5 C_d + 3,
\]
such that every $F \in \Fd$ satisfying
\[
\inf_{\abs{w}=1}\norm{wF - 1}_{\Fd}
\le
\delta_d
\]
also satisfies
\[
\inf_{\abs{w}=1}\norm{wF - 1}_{\Fd}
\le
\widetilde M_d\, \bigl\|\, \abs{F} - 1 \,\bigr\|_{\Fd}.
\]
\end{corollary}

\begin{proposition}[Optional soft global consequence]\label{prop:soft-global}
Fix $d \ge 1$.  There exists a constant $\widehat M_d > 0$ such that every
$F \in \Fd$ satisfies
\[
\inf_{\abs{w}=1}\norm{wF - 1}_{\Fd}
\le
\widehat M_d\, \bigl\|\, \abs{F} - 1 \,\bigr\|_{\Fd}.
\]
\end{proposition}

Theorem~\ref{thm:orthog} implies Corollary~\ref{cor:small} by the
orthogonal reduction recorded in Proposition~\ref{prop:orth-reduction}.
Proposition~\ref{prop:soft-global} then follows from
Corollary~\ref{cor:small} by a separate compactness argument, once phase
retrieval at the constant is verified.  These auxiliary steps are carried
out in Section~\ref{sec:prelim-local}.  They are included for completeness,
but the quantitative core of the paper stops at
Corollary~\ref{cor:small}.

We first prove Theorem~\ref{thm:orthog} for holomorphic polynomials with
vanishing constant term.  The passage to general $U \in \Fd$ is by
total-degree truncation and is completed in Section~\ref{sec:assembly}.

\section{Fock-Space Preliminaries and Optional Soft Reduction}\label{sec:prelim-local}

We record the standard several-variable Fock-space facts used later to
identify $1^\perp$ with the functions vanishing at the origin and to pass
from polynomial estimates to general elements of $\Fd$.

For a multi-index $\alpha = (\alpha_1,\dots,\alpha_d) \in \N^d$, write
\[
\abs{\alpha} := \alpha_1 + \cdots + \alpha_d,
\qquad
\alpha! := \alpha_1! \cdots \alpha_d!,
\qquad
z^\alpha := z_1^{\alpha_1}\cdots z_d^{\alpha_d}.
\]
For $a \in \C^d$, set
\[
K_a(z) := \exp\!\Bigl(\sum_{j=1}^d z_j \overline{a_j}\Bigr).
\]

\begin{proposition}[Kernel, basis, and truncations]\label{prop:fock-facts}
For $\alpha \in \N^d$, define
\[
e_\alpha(z) := \frac{z^\alpha}{\sqrt{\alpha!}}.
\]
Then:
\begin{enumerate}[(1)]
\item $\{e_\alpha : \alpha \in \N^d\}$ is an orthonormal basis of $\Fd$.
\item For every $a \in \C^d$, the function $K_a$ belongs to $\Fd$ and
\[
\langle F, K_a\rangle_{\Fd} = F(a)
\qquad
\text{for every } F \in \Fd.
\]
\item In particular, $K_0 = 1$, so
\[
\langle F,1\rangle_{\Fd} = F(0)
\qquad
\text{for every } F \in \Fd.
\]
\item If
\[
F = \sum_{\alpha \in \N^d} c_\alpha e_\alpha,
\qquad
F^{(\le N)} := \sum_{\abs{\alpha}\le N} c_\alpha e_\alpha,
\]
then $F^{(\le N)} \to F$ in $\Fd$ as $N \to \infty$.
\end{enumerate}
\end{proposition}

\begin{proof}
Orthogonality is a product of the one-variable Gaussian moments:
\[
\langle e_\alpha,e_\beta\rangle_{\Fd}
=
\prod_{j=1}^d
\frac{1}{\pi}\int_{\C}
\frac{z_j^{\alpha_j}\overline{z_j}^{\beta_j}}{\sqrt{\alpha_j!\beta_j!}}
\ee^{-\abs{z_j}^2}\,\dd m(z_j)
=
\delta_{\alpha\beta}.
\]

Now let $F \in \Fd$, and write its power-series expansion
\[
F(z) = \sum_{\alpha \in \N^d} b_\alpha z^\alpha.
\]
Because $F$ is entire, this series converges absolutely and uniformly on
every compact subset of $\C^d$.  Fix
$r=(r_1,\dots,r_d)\in(0,\infty)^d$.  On the polytorus
\[
\bigl\{(r_1 \ee^{\ii\theta_1},\dots,r_d \ee^{\ii\theta_d}) :
(\theta_1,\dots,\theta_d)\in[0,2\pi]^d\bigr\}
\]
we may therefore apply Parseval on $\mathbb T^d$ to obtain
\[
\int_{[0,2\pi]^d}
\abs{F(r_1 \ee^{\ii\theta_1},\dots,r_d \ee^{\ii\theta_d})}^2\,
\prod_{j=1}^d \frac{\dd\theta_j}{2\pi}
=
\sum_{\alpha \in \N^d} \abs{b_\alpha}^2 r^{2\alpha},
\]
where $r^{2\alpha}:=r_1^{2\alpha_1}\cdots r_d^{2\alpha_d}$.  Integrating in
$r$ against $\prod_{j=1}^d 2 r_j \ee^{-r_j^2}\,\dd r_j$ and using Tonelli's
theorem gives
\[
\norm{F}_{\Fd}^2
=
\sum_{\alpha \in \N^d} \abs{b_\alpha}^2
\prod_{j=1}^d
\left(
2\int_0^\infty r_j^{2\alpha_j+1}\ee^{-r_j^2}\,\dd r_j
\right)
=
\sum_{\alpha \in \N^d} \abs{b_\alpha}^2 \alpha!.
\]
Thus $c_\alpha := b_\alpha \sqrt{\alpha!}$ belongs to $\ell^2(\N^d)$ and
\[
F
=
\sum_{\alpha \in \N^d} c_\alpha e_\alpha,
\qquad
\norm{F}_{\Fd}^2
=
\sum_{\alpha \in \N^d} \abs{c_\alpha}^2.
\]
In particular, $\{e_\alpha\}$ is complete as well as orthonormal, hence an
orthonormal basis.  For the total-degree truncations,
\[
F^{(\le N)} := \sum_{\abs{\alpha}\le N} c_\alpha e_\alpha,
\]
we therefore have
\[
\norm{F - F^{(\le N)}}_{\Fd}^2
=
\sum_{\abs{\alpha}>N} \abs{c_\alpha}^2
\xrightarrow[N\to\infty]{} 0.
\]

For $a \in \C^d$,
\[
\abs{K_a(z)}^2 \ee^{-\abs{z}^2}
=
\ee^{2\operatorname{Re}\sum_j z_j\overline{a_j} - \abs{z}^2}
=
\ee^{\abs{a}^2}\,\ee^{-\abs{z-a}^2},
\]
so $K_a \in \Fd$.  Moreover,
\[
K_a(z)
=
\sum_{\alpha \in \N^d} \frac{\overline{a}^\alpha}{\alpha!}\, z^\alpha
=
\sum_{\alpha \in \N^d} \frac{\overline{a}^\alpha}{\sqrt{\alpha!}}\, e_\alpha(z),
\]
and the coefficient sequence is square-summable because
\[
\sum_{\alpha \in \N^d} \frac{\abs{a}^{2\alpha}}{\alpha!}
=
\prod_{j=1}^d \sum_{n\ge0}\frac{\abs{a_j}^{2n}}{n!}
=
\ee^{\abs{a}^2}.
\]
\[
\norm{K_a}_{\Fd}^2
=
\sum_{\alpha \in \N^d} \frac{\abs{a}^{2\alpha}}{\alpha!}
=
\ee^{\abs{a}^2},
\]
in agreement with the direct Gaussian computation above.  If
$F=\sum_\alpha c_\alpha e_\alpha$, then
Cauchy--Schwarz gives absolute convergence of
\[
\sum_{\alpha \in \N^d} c_\alpha \frac{a^\alpha}{\sqrt{\alpha!}},
\]
because
\[
\sum_{\alpha \in \N^d} \frac{\abs{a}^{2\alpha}}{\alpha!}
=
\ee^{\abs{a}^2}
<
\infty.
\]
Therefore
\[
F(a)
=
\sum_{\alpha \in \N^d} c_\alpha \frac{a^\alpha}{\sqrt{\alpha!}}
=
\left\langle
\sum_{\alpha \in \N^d} c_\alpha e_\alpha,\,
\sum_{\alpha \in \N^d} \frac{\overline{a}^\alpha}{\sqrt{\alpha!}}\,e_\alpha
\right\rangle_{\Fd}
=
\langle F, K_a\rangle_{\Fd}.
\]
Taking $a=0$ gives $K_0=1$ and $\langle F,1\rangle_{\Fd}=F(0)$.
\end{proof}

\begin{proposition}[Orthogonal reduction at the constant]\label{prop:orth-reduction}
Let $C_d$ be as in Theorem~\ref{thm:orthog}, and set
\[
\delta_d := \frac{1}{C_d+1},
\qquad
\widetilde M_d := 5C_d+3.
\]
Suppose $h \in \Fd$ satisfies
\[
\norm{h}_{\Fd} \le \delta_d
\qquad\text{and}\qquad
\langle h,1\rangle_{\Fd} \in \mathbb R.
\]
Then
\[
\norm{h}_{\Fd}
\le
\widetilde M_d\,\norm{\abs{1+h}-1}_{\Fd}.
\]
\end{proposition}

\begin{proof}
This is the abstract Hilbert-space orthogonal reduction from
\texttt{Constant/Latex/FockSPR/OrthogonalReduction.tex}, applied with
$H=\Fd$, $f_0=1$, and $M=C_d$.  By
Proposition~\ref{prop:fock-facts}, the orthogonality condition $g \perp 1$
is equivalent to $g(0)=0$, so Theorem~\ref{thm:orthog} supplies the
required orthogonal coercivity hypothesis.  The quoted reduction theorem
then yields the stated constants $\delta_d = 1/(C_d+1)$ and
$\widetilde M_d = 5C_d+3$.
\end{proof}

\begin{proof}[Proof of Corollary~\ref{cor:small}]
Let $F \in \Fd$ satisfy
\[
\inf_{\abs{w}=1}\norm{wF-1}_{\Fd} \le \delta_d.
\]
Because the unit circle is compact and
$w \mapsto \norm{wF-1}_{\Fd}$ is continuous, there exists a minimising
phase $w_0 \in \C$, $\abs{w_0}=1$, such that
\[
\norm{w_0F-1}_{\Fd}
=
\inf_{\abs{w}=1}\norm{wF-1}_{\Fd}
\le
\delta_d.
\]
For every unimodular $w$,
\[
\norm{wF-1}_{\Fd}^2
=
\norm{F}_{\Fd}^2 + 1 - 2\operatorname{Re}\!\bigl(w\langle F,1\rangle_{\Fd}\bigr),
\]
because $\norm{wF}_{\Fd}=\norm{F}_{\Fd}$ and the inner product is linear in
the first slot.  Since $w_0$ minimises the left-hand side, it maximises
$\operatorname{Re}(w\langle F,1\rangle_{\Fd})$ over $\abs{w}=1$.  If
$\langle F,1\rangle_{\Fd}=0$, this quantity is identically zero and the next
display is trivial.  Otherwise its maximum is
$\abs{\langle F,1\rangle_{\Fd}}$, attained exactly when
\[
w_0 \langle F,1\rangle_{\Fd}
=
\bigl|\langle F,1\rangle_{\Fd}\bigr|
\in
\mathbb R_{\ge 0},
\]
so $h := w_0F-1$ satisfies $\langle h,1\rangle_{\Fd} \in \mathbb R$.
Since $\norm{h}_{\Fd}\le \delta_d$, Proposition~\ref{prop:orth-reduction}
therefore yields
\[
\norm{w_0F-1}_{\Fd}
\le
\widetilde M_d\,\bigl\|\, \abs{F} - 1 \,\bigr\|_{\Fd}.
\]
Since the left-hand side is the phase distance to $1$, this is precisely
the claimed estimate.
\end{proof}

\begin{proposition}[Phase retrieval at the constant]\label{prop:phase-at-one}
If $F \in \Fd$ satisfies $\abs{F}=1$ almost everywhere on $\C^d$, then
$F \equiv w$ for some $w \in \C$ with $\abs{w}=1$.
\end{proposition}

\begin{proof}
Because the Gaussian density $\ee^{-\abs{z}^2}$ never vanishes,
Gaussian-null sets and Lebesgue-null sets coincide.  Since $F$ is
continuous, the almost-everywhere identity $\abs{F}=1$ forces
$\abs{F(z)}=1$ for every $z \in \C^d$.

Fix $(z_2,\dots,z_d) \in \C^{d-1}$.  Then the function
\[
t \mapsto F(t,z_2,\dots,z_d)
\]
is entire on $\C$ and bounded by $1$, hence constant by Liouville's
theorem.  Thus $F$ is independent of $z_1$.  Repeating the same argument in
each coordinate shows that $F$ is constant on $\C^d$.  Its constant value
has modulus~$1$.
\end{proof}

\begin{remark}[On the optional soft local-to-global step]
Proposition~\ref{prop:local-reduction} is the only non-quantitative part of
the argument: it uses weak compactness, pointwise continuity of evaluation,
and subsequence extraction from $L^2$ convergence.  The quantitative core of
the note is Theorem~\ref{thm:orthog} together with
Corollary~\ref{cor:small}; Proposition~\ref{prop:local-reduction} is only
used to pass from that small-data statement to the global estimate recorded
as Proposition~\ref{prop:soft-global}.
\end{remark}

\begin{proposition}[Local reduction at the constant]\label{prop:local-reduction}
Suppose that there exist $\varepsilon > 0$ and $C_{\mathrm{loc}} > 0$ such
that every $F \in \Fd$ satisfying
\[
\inf_{\abs{w}=1}\norm{wF - 1}_{\Fd} < \varepsilon
\]
also satisfies
\[
\inf_{\abs{w}=1}\norm{wF - 1}_{\Fd}
\le
C_{\mathrm{loc}}\,\bigl\|\, \abs{F} - 1 \,\bigr\|_{\Fd}.
\]
Then there exists $C_{\mathrm{glob}} > 0$ such that
\[
\inf_{\abs{w}=1}\norm{wF - 1}_{\Fd}
\le
C_{\mathrm{glob}}\,\bigl\|\, \abs{F} - 1 \,\bigr\|_{\Fd}
\qquad
\text{for every } F \in \Fd.
\]
\end{proposition}

\begin{proof}
Assume the contrary.  Then for each $n \ge 1$ there exists $F_n \in \Fd$
such that
\[
D_n > nR_n,
\]
where
\[
D_n := \inf_{\abs{w}=1}\norm{wF_n - 1}_{\Fd},
\qquad
R_n := \bigl\|\, \abs{F_n} - 1 \,\bigr\|_{\Fd}.
\]

If $\norm{F_n}_{\Fd} \ge 4$, then
\[
D_n \le \norm{F_n}_{\Fd} + 1 \le \frac54 \norm{F_n}_{\Fd},
\]
whereas
\[
R_n
\ge
\bigl|\norm{F_n}_{\Fd} - \norm{1}_{\Fd}\bigr|
=
\norm{F_n}_{\Fd} - 1
\ge
\frac34 \norm{F_n}_{\Fd}.
\]
Hence $D_n \le \frac53 R_n$, contradicting $D_n > nR_n$ for every
$n \ge 2$.  After discarding finitely many terms, we may therefore assume
that $\norm{F_n}_{\Fd} < 4$ for all $n$.  In particular, $D_n \le 5$, so
$R_n < 5/n \to 0$.

The sequence $(F_n)$ is bounded in the Hilbert space $\Fd$, so after
passing to a subsequence we may assume that $F_n \rightharpoonup F$ weakly
in $\Fd$ for some $F \in \Fd$.  By
Proposition~\ref{prop:fock-facts}, point evaluations are continuous on
$\Fd$, so $F_n(z) \to F(z)$ for every $z \in \C^d$.

Since $R_n \to 0$ in $L^2$, after passing to a further subsequence we may
also assume that $\abs{F_n(z)} \to 1$ for almost every $z$.  Combining this
with the pointwise convergence $F_n(z) \to F(z)$ yields
$\abs{F(z)} = 1$ for almost every $z$.  Proposition~\ref{prop:phase-at-one}
therefore implies that $F \equiv w_\ast$ for some unimodular constant
$w_\ast$.

Moreover,
\[
\bigl|\norm{F_n}_{\Fd} - 1\bigr|
\le
R_n
\xrightarrow[n\to\infty]{} 0,
\]
so $\norm{F_n}_{\Fd} \to 1 = \norm{F}_{\Fd}$.  Weak convergence together
with convergence of norms now gives $F_n \to F$ strongly in $\Fd$.  Hence
\[
D_n
\le
\norm{\overline{w_\ast}F_n - 1}_{\Fd}
\xrightarrow[n\to\infty]{} 0.
\]
For all sufficiently large $n$ we therefore have $D_n < \varepsilon$, so
the local hypothesis gives $D_n \le C_{\mathrm{loc}} R_n$.  This
contradicts $D_n > nR_n$ once $n > C_{\mathrm{loc}}$.
\end{proof}

\begin{proof}[Proof of Proposition~\ref{prop:soft-global}]
Apply Proposition~\ref{prop:local-reduction} with
$\varepsilon = \delta_d$ and $C_{\mathrm{loc}} = \widetilde M_d$,
using Corollary~\ref{cor:small}.  The phase-retrieval input needed in the
proof of Proposition~\ref{prop:local-reduction} is exactly
Proposition~\ref{prop:phase-at-one}.  The resulting constant
$C_{\mathrm{glob}}$ is the desired $\widehat M_d$.
\end{proof}

\section{Polar Coordinates and Homogeneous Pieces}\label{sec:polar}

Let $\sigma_d$ denote the probability measure on the unit sphere
$\Sphere \subset \C^d$.  Since
\[
\abs{\Sphere}
=
\frac{2\pi^d}{(d-1)!},
\]
polar coordinates give
\begin{equation}\label{eq:polar}
\norm{F}_{\Fd}^2
=
\frac{2}{(d-1)!}
\int_0^\infty
r^{2d-1} \ee^{-r^2}
\biggl(
\int_{\Sphere} \abs{F(r\omega)}^2\,\dd \sigma_d(\omega)
\biggr)\dd r.
\end{equation}

For the polynomial part of the proof, let $U$ be a holomorphic polynomial
with $U(0)=0$.  Write
\[
U = \sum_{n=1}^D H_n,
\]
where each $H_n$ has total degree $n$.

\begin{lemma}[Sphere orthogonality by degree]\label{lem:sphere-orth}
If $n \ne m$, then
\[
\int_{\Sphere}
H_n(\omega)\, \overline{H_m(\omega)}\, \dd \sigma_d(\omega)
=
0.
\]
\end{lemma}

\begin{proof}
The measure $\sigma_d$ is invariant under the global phase action
$\omega \mapsto \ee^{\ii t}\omega$.  Since
$H_n(\ee^{\ii t}\omega) = \ee^{\ii nt} H_n(\omega)$ and similarly for $H_m$,
the integral is multiplied by $\ee^{\ii(n-m)t}$ for every $t$, hence must
vanish.
\end{proof}

Consequently, on each sphere of radius $r$,
\[
\int_{\Sphere} \abs{U(r\omega)}^2\,\dd \sigma_d(\omega)
=
\sum_{n=1}^D r^{2n}\, \norm{H_n}_{L^2(\Sphere)}^2,
\]
\[
\norm{U}_{\Fd}^2
=
\frac{2}{(d-1)!}
\sum_{n=1}^D
\norm{H_n}_{L^2(\Sphere)}^2
\int_0^\infty r^{2n+2d-1}\ee^{-r^2}\,\dd r.
\]
With the substitution $s=r^2$,
\[
2\int_0^\infty r^{2n+2d-1}\ee^{-r^2}\,\dd r
=
\int_0^\infty s^{n+d-1}\ee^{-s}\,\dd s
=
(n+d-1)!,
\]
and therefore
\begin{equation}\label{eq:fock-by-degree}
\norm{U}_{\Fd}^2
=
\sum_{n=1}^D
\frac{(n+d-1)!}{(d-1)!}\,
\norm{H_n}_{L^2(\Sphere)}^2.
\end{equation}

\section{Sphere Versions of the Circle Estimates}

The one-dimensional proof uses two circle estimates:

\begin{itemize}
\item the local circle estimate
$\norm{P}_{L^2(S^1)}^2 \le 144 L\, \norm{\rhoFn \circ P}_{L^2(S^1)}^2$
for trigonometric polynomials with $L$ positive frequencies;
\item the high-frequency estimate
$\norm{P}_{L^2(S^1)}^2 \le 32\, \norm{\rhoFn \circ P}_{L^2(S^1)}^2$
for bands $\{N,\dots,N+L-1\}$ with $1343L^2 \le N^2$.
\end{itemize}

Here $\sigma$ denotes the probability Haar measure on $S^1$.

The point is that both lift directly to $\Sphere$.

\begin{proposition}[Lifted local estimate]\label{prop:lifted-local}
Let $E \subset \N_{\ge 1}$ be finite with $\abs{E} = L$, and let
\[
P(\omega) = \sum_{n \in E} G_n(\omega),
\]
where each $G_n$ is the restriction to $\Sphere$ of a homogeneous
holomorphic polynomial of degree $n$.  Then
\[
\norm{P}_{L^2(\Sphere)}^2
\le
144 L\, \norm{\rhoFn \circ P}_{L^2(\Sphere)}^2.
\]
\end{proposition}

\begin{proof}
Fix $\omega \in \Sphere$.  The function
\[
t \mapsto P(\ee^{\ii t}\omega)
=
\sum_{n \in E} \ee^{\ii nt} G_n(\omega)
\]
is a trigonometric polynomial with the same positive frequency set $E$.
The one-dimensional local circle estimate applies to this polynomial
uniformly in the coefficients $G_n(\omega)$, so
\[
\int_{S^1} \abs{P(\ee^{\ii t}\omega)}^2\,\dd \sigma(t)
\le
144 L
\int_{S^1} \rhoFn(P(\ee^{\ii t}\omega))^2\,\dd \sigma(t).
\]
Integrating in $\omega$ and using Tonelli's theorem gives
\[
\int_{\Sphere}\!\int_{S^1} \abs{P(\ee^{\ii t}\omega)}^2
\,\dd \sigma(t)\,\dd \sigma_d(\omega)
\le
144L
\int_{\Sphere}\!\int_{S^1} \rhoFn(P(\ee^{\ii t}\omega))^2
\,\dd \sigma(t)\,\dd \sigma_d(\omega).
\]
For each fixed $t$, the map $\omega \mapsto \ee^{\ii t}\omega$ preserves
$\sigma_d$, hence
\[
\int_{\Sphere} \abs{P(\ee^{\ii t}\omega)}^2\,\dd \sigma_d(\omega)
=
\int_{\Sphere} \abs{P(\omega)}^2\,\dd \sigma_d(\omega),
\]
and likewise
\[
\int_{\Sphere} \rhoFn(P(\ee^{\ii t}\omega))^2\,\dd \sigma_d(\omega)
=
\int_{\Sphere} \rhoFn(P(\omega))^2\,\dd \sigma_d(\omega).
\]
Substituting these identities and integrating in $t$ yields the claimed
estimate.
\end{proof}

\begin{proposition}[Lifted high-frequency estimate]\label{prop:lifted-high}
Let $N,L \ge 1$ satisfy $1343L^2 \le N^2$, and let
\[
P(\omega) = \sum_{m=0}^{L-1} G_{N+m}(\omega),
\]
where each $G_{N+m}$ is homogeneous holomorphic of degree $N+m$.  Then
\[
\norm{P}_{L^2(\Sphere)}^2
\le
32\, \norm{\rhoFn \circ P}_{L^2(\Sphere)}^2.
\]
\end{proposition}

\begin{proof}
For fixed $\omega$, the function
$t \mapsto P(\ee^{\ii t}\omega)$ has frequencies contained in
$\{N,\dots,N+L-1\}$, with zero coefficients allowed.  The one-dimensional
high-frequency estimate therefore applies uniformly in $\omega$, yielding
\[
\int_{S^1} \abs{P(\ee^{\ii t}\omega)}^2\,\dd \sigma(t)
\le
32\int_{S^1} \rhoFn(P(\ee^{\ii t}\omega))^2\,\dd \sigma(t).
\]
Now integrate in $\omega$ and use the same Tonelli-plus-invariance argument
as in Proposition~\ref{prop:lifted-local}.
\end{proof}

\section{Degree Blocks and Leakage}

The one-dimensional proof groups monomials by the radius where their Gaussian
mass is largest.  In $\C^d$, the homogeneous degree-$n$ piece plays the same
role, because its radial weight is
\[
r^{2n+2d-1}\ee^{-r^2}.
\]
Set
\[
m_n := n+d-1,
\qquad
r_n := \sqrt{m_n + \tfrac12} = \sqrt{n+d-\tfrac12}.
\]
Then $r_n$ is the unique maximiser of the function
\[
\phi_n(r)
:=
(2m_n+1)\log r - r^2.
\]

Fix an integer window width $M \ge 2$, to be chosen later.
For $\ell \ge 1$, define the degree block
\[
I_\ell^{(d)}
:=
\bigl\{
n \in \N_{\ge 1} : \ell^2 \le m_n < (\ell+1)^2
\bigr\},
\]
and the corresponding block polynomial
\[
U_\ell := \sum_{n \in I_\ell^{(d)}} H_n.
\]

\begin{lemma}[Peak localisation]\label{lem:peak}
If $n \in I_\ell^{(d)}$, then $\ell < r_n < \ell+1$.
\end{lemma}

\begin{proof}
Since $m_n$ is an integer and $m_n < (\ell+1)^2$, we in fact have
$m_n \le (\ell+1)^2 - 1$.  Therefore
\[
\ell^2
<
m_n + \tfrac12
\le
(\ell+1)^2 - \tfrac12
<
(\ell+1)^2,
\]
which is equivalent to $\ell < r_n < \ell+1$.
\end{proof}

For each annulus index $j \ge 0$, define
\[
V_j := \sum_{\abs{\ell-j}\le M} U_\ell,
\qquad
R_j := U - V_j.
\]
Thus $V_j$ contains the degrees whose radial peaks lie within $M$ annuli of
$[j,j+1]$, and $R_j$ is the distant remainder.

\begin{lemma}[Degree bands on an annulus]\label{lem:count}
For every $j \ge 0$, set
\[
a_j := \max(1,j-M),
\qquad
b_j := j+M,
\qquad
N_j := \max\bigl(1,\,a_j^2 - d + 1\bigr),
\qquad
K_j := (b_j+1)^2 - d.
\]
Then every degree appearing in $V_j$ lies in the possibly empty interval
\[
\bigl\{
n \in \N_{\ge 1} : N_j \le n \le K_j
\bigr\}.
\]
In particular, this interval has length at most
\[
L_j := (b_j+1)^2 - a_j^2
\]
and therefore $V_j$ uses at most $L_j$ distinct degrees.
If moreover $j \ge M+1$, then $a_j = j-M$ and
\[
L_j = (2M+1)(2j+1).
\]
\end{lemma}

\begin{proof}
If $\abs{\ell-j}\le M$, then $a_j \le \ell \le b_j$.  Hence every
$n \in I_\ell^{(d)}$ satisfies
\[
a_j^2 \le \ell^2 \le m_n < (\ell+1)^2 \le (b_j+1)^2.
\]
Since $m_n = n+d-1$, this implies
\[
a_j^2-d+1 \le n \le (b_j+1)^2-d.
\]
Together with $n \ge 1$, this shows that every degree appearing in $V_j$
lies in the stated interval, and that interval has length
\[
K_j-N_j+1 \le (b_j+1)^2-a_j^2 = L_j.
\]
If $j \ge M+1$, then $a_j = j-M$, and the displayed formula simplifies to
\[
(j+M+1)^2-(j-M)^2 = (2M+1)(2j+1).
\]
\end{proof}

The radial saddle-point argument follows the exact Gamma-integral proof from
\texttt{Constant/Latex/FockSPR/BlockDecomposition.tex}, with $n$ replaced by
$m_n$.

\begin{lemma}[Single-degree radial bound]\label{lem:single-degree}
For every $n \ge 1$ and every $r > 0$,
\[
\ee^{\phi_n(r)}
\le
\ee^{\phi_n(r_n)} \ee^{-(r-r_n)^2},
\qquad
\ee^{\phi_n(r_n)}
\le
\frac{\ee^{1/4}}{2}\, m_n!.
\]
\end{lemma}

\begin{proof}
Since
\[
\phi_n'(r_n)=0,
\qquad
\phi_n''(r)
=
-\frac{2m_n+1}{r^2}-2
\le
-2,
\]
Taylor's theorem gives
\[
\phi_n(r)
\le
\phi_n(r_n)-(r-r_n)^2
\qquad (r>0),
\]
which is the first estimate.

For the second estimate, note that $m_n \ge d \ge 1$, so
\[
\frac{m_n!}{2}
=
\int_0^\infty \ee^{\phi_n(r)}\,\dd r.
\]
Since $m_n \ge 1$, we have $r_n = \sqrt{m_n+\tfrac12} \ge \sqrt{3/2} > 1$,
so the interval $[r_n-\tfrac12,r_n+\tfrac12]$ lies in $(0,\infty)$.  By the
first estimate,
\[
\phi_n(r) \ge \phi_n(r_n)-\frac14
\qquad
\text{for every } r \in [r_n-\tfrac12,r_n+\tfrac12].
\]
Therefore
\[
\frac{m_n!}{2}
\ge
\int_{r_n-\frac12}^{r_n+\frac12} \ee^{\phi_n(r)}\,\dd r
\ge
\ee^{\phi_n(r_n)-\frac14}.
\]
Rearranging gives
\[
\ee^{\phi_n(r_n)}
\le
\frac{\ee^{1/4}}{2}\,m_n!,
\]
as claimed.
\end{proof}

\begin{proposition}[Single-block leakage]\label{prop:single-leak}
For every $\ell \ge 1$ and $j \ge 0$,
\[
\Ann{j}(U_\ell)
\le
\ee^{1/4}\,
\ee^{-(\abs{j-\ell}-1)_+^2}
\norm{U_\ell}_{\Fd}^2,
\]
where
\[
\Ann{j}(F)
:=
\frac{2}{(d-1)!}
\int_j^{j+1} r^{2d-1}\ee^{-r^2}
\biggl(
\int_{\Sphere} \abs{F(r\omega)}^2\,\dd \sigma_d(\omega)
\biggr)\dd r.
\]
\end{proposition}

\begin{proof}
Because distinct degrees are orthogonal on $\Sphere$,
\[
\int_{\Sphere}\abs{U_\ell(r\omega)}^2\,\dd \sigma_d(\omega)
=
\sum_{n \in I_\ell^{(d)}} r^{2n}\norm{H_n}_{L^2(\Sphere)}^2.
\]
Hence
\[
\Ann{j}(U_\ell)
=
\frac{2}{(d-1)!}
\sum_{n \in I_\ell^{(d)}} \norm{H_n}_{L^2(\Sphere)}^2
\int_j^{j+1} r^{2m_n+1}\ee^{-r^2}\,\dd r.
\]
Lemma~\ref{lem:single-degree} gives
\[
\int_j^{j+1} r^{2m_n+1}\ee^{-r^2}\,\dd r
\le
\frac{\ee^{1/4}}{2}\,m_n!\,
\ee^{-\dist(r_n,[j,j+1])^2}.
\]
If $n \in I_\ell^{(d)}$, then Lemma~\ref{lem:peak} yields
$r_n \in (\ell,\ell+1)$, and therefore
\[
\dist(r_n,[j,j+1]) \ge (\abs{j-\ell}-1)_+.
\]
Substituting this into the previous display yields
\[
\Ann{j}(U_\ell)
\le
\frac{\ee^{1/4}}{(d-1)!}\,
\ee^{-(\abs{j-\ell}-1)_+^2}
\sum_{n \in I_\ell^{(d)}} m_n!\,\norm{H_n}_{L^2(\Sphere)}^2.
\]
The sum is exactly $(d-1)!\,\norm{U_\ell}_{\Fd}^2$ by
\eqref{eq:fock-by-degree}, so the stated bound follows.
\end{proof}

\begin{corollary}[Total leakage]\label{cor:leakage}
Let
\[
\eta_M := 2\ee^{1/4}\sum_{m \ge M}\ee^{-m^2}.
\]
Then
\[
\sum_{j \ge 0} \Ann{j}(R_j)
\le
\eta_M\, \norm{U}_{\Fd}^2.
\]
\end{corollary}

\begin{proof}
On each annulus, the blocks $U_\ell$ with $\abs{\ell-j}>M$ have pairwise
disjoint degree supports.  By Lemma~\ref{lem:sphere-orth}, this gives
\[
\Ann{j}(R_j)
=
\sum_{\abs{\ell-j}>M} \Ann{j}(U_\ell).
\]
Applying Proposition~\ref{prop:single-leak} termwise yields
\[
\Ann{j}(R_j)
\le
\ee^{1/4}\sum_{\abs{\ell-j}>M}
\ee^{-(\abs{j-\ell}-1)_+^2}\norm{U_\ell}_{\Fd}^2.
\]
Summing over $j \ge 0$ and using Tonelli's theorem,
\[
\sum_{j \ge 0} \Ann{j}(R_j)
\le
\ee^{1/4}\sum_{\ell \ge 1}\norm{U_\ell}_{\Fd}^2
\sum_{\substack{j \ge 0\\ \abs{\ell-j}>M}}
\ee^{-(\abs{j-\ell}-1)_+^2}.
\]
If $\abs{\ell-j}>M$ and $M \ge 2$, then
$(\abs{j-\ell}-1)_+ = \abs{j-\ell}-1 \ge M$.  For each integer $m \ge M$,
there are at most two indices $j \ge 0$ such that $\abs{j-\ell}-1 = m$,
so
\[
\sum_{\substack{j \ge 0\\ \abs{\ell-j}>M}}
\ee^{-(\abs{j-\ell}-1)_+^2}
\le
2\sum_{m \ge M}\ee^{-m^2}.
\]
Finally, the blocks are orthogonal in $\Fd$ by \eqref{eq:fock-by-degree}, so
\[
\sum_{\ell \ge 1}\norm{U_\ell}_{\Fd}^2
=
\norm{U}_{\Fd}^2.
\]
Combining the previous displays gives the stated bound.
\end{proof}

\section{Annulus Estimate}

For fixed $j$ and $r \in [j,j+1]$, the restriction
$\omega \mapsto V_j(r\omega)$ uses only positive degrees.  By
Lemma~\ref{lem:count}, these degrees lie in a contiguous band starting at
degree at least $N_j$ and having length at most $L_j$.

Choose an integer $J = J(d,M) \ge M+1$ so large that
\begin{equation}\label{eq:threshold}
1343\, L_j^2 \le N_j^2
\qquad
\text{for all } j \ge J.
\end{equation}
Such a $J$ exists because $L_j = O(j)$ whereas $N_j \sim j^2$.
For instance, one may take
\[
J(d,M) := M + 75(2M+1) + \lceil \sqrt d \rceil,
\]
because for $j = M + 75(2M+1) + s$ with $s \ge \sqrt d$,
\[
N_j \ge (75(2M+1)+s)^2 - d + 1
\ge 5625(2M+1)^2 + 150(2M+1)s + 1,
\]
whereas
\[
\sqrt{1343}\,L_j
\le
37(2M+1)(2j+1)
=
37(2M+1)\bigl(151(2M+1)+2s\bigr).
\]
Since $5625 > 37 \cdot 151$ and $150 > 74$, the required inequality holds.

\begin{remark}[Illustrative threshold when $M=5$]\label{rem:explicit-J}
If one fixes $M=5$, then one may for instance take
\[
J(d,5) = 819 + \lceil \sqrt d \rceil.
\]
Indeed, if $j = 819+s$ with $s \ge \sqrt d$, then
\[
N_j \ge (j-5)^2 - d + 1 = (814+s)^2 - d + 1 \ge 662597 + 1628s,
\]
while
\[
\sqrt{1343}\,L_j
\le
11\sqrt{1343}\,(2j+1)
<
404(2j+1)
=
662156 + 808s.
\]
This special case is included only to compare with the one-dimensional
bookkeeping.  The proof of Theorem~\ref{thm:orthog} below does not fix
$M=5$; instead, it later chooses $M$ depending on $d$ so that the absorption
step closes.
\end{remark}

\begin{proposition}[Uniform annulus estimate]\label{prop:annulus}
Set
\[
B_{d,M}^2 := 144(2M+1)(2J+1).
\]
Then for every $j \ge 0$ and every $r \in [j,j+1]$,
\[
\int_{\Sphere} \abs{V_j(r\omega)}^2\,\dd \sigma_d(\omega)
\le
B_{d,M}^2
\int_{\Sphere} \rhoFn(V_j(r\omega))^2\,\dd \sigma_d(\omega).
\]
\end{proposition}

\begin{proof}
If $V_j=0$, there is nothing to prove.  Assume henceforth that $V_j \ne 0$.

If $j < J$, apply Proposition~\ref{prop:lifted-local} together with the
bound $L_j \le (2M+1)(2J+1)$.  Indeed, if $j \ge M+1$, then
Lemma~\ref{lem:count} gives
\[
L_j = (2M+1)(2j+1) \le (2M+1)(2J+1).
\]
If instead $0 \le j \le M$, then
\[
L_j = (j+M+1)^2-1 \le (2M+1)^2-1 \le (2M+1)(2J+1),
\]
because $J \ge M+1$.

If $j \ge J$, let
\[
L'_j := K_j-N_j+1.
\]
Lemma~\ref{lem:count} shows that the degrees of $V_j$ lie in the
contiguous band $\{N_j,\dots,K_j\}$, and $L'_j \le L_j$.  Padding absent
degrees in this band by zero, we may write
\[
V_j(r\omega)
=
\sum_{m=0}^{L'_j-1} \widetilde G_{N_j+m}(\omega),
\]
where each $\widetilde G_{N_j+m}$ is homogeneous holomorphic of degree
$N_j+m$ and may be identically zero.  Proposition~\ref{prop:lifted-high}
therefore applies with bandwidth $L'_j$.  Since
\eqref{eq:threshold} gives
\[
1343(L'_j)^2 \le 1343L_j^2 \le N_j^2,
\]
this yields the stronger estimate
\[
\int_{\Sphere} \abs{V_j(r\omega)}^2\,\dd \sigma_d(\omega)
\le
32
\int_{\Sphere} \rhoFn(V_j(r\omega))^2\,\dd \sigma_d(\omega),
\]
and $32 \le B_{d,M}^2$.
\end{proof}

\section{Assembly}\label{sec:assembly}

We now assemble the annulus estimate and the leakage bound exactly as in the
one-dimensional proof, with the higher-dimensional ingredients proved above
replacing their circle counterparts.

\begin{proof}[Proof of Theorem~\ref{thm:orthog}]
We first treat the case where $U$ is a holomorphic polynomial with
$U(0)=0$.

Fix $d$.  Since $\eta_M \to 0$ super-exponentially as $M \to \infty$,
while the explicit choice of $J(d,M)$ above gives $B_{d,M} = O_d(M)$,
we may choose $M \ge 2$ so large that
\[
(4B_{d,M}^2+2)\eta_M < \frac12.
\]
Now choose this $M$, set $J = J(d,M)$, and let $B = B_{d,M}$.

By Proposition~\ref{prop:annulus}, for every $r \in [j,j+1]$,
\[
\int_{\Sphere} \abs{V_j(r\omega)}^2\,\dd \sigma_d(\omega)
\le
B^2
\int_{\Sphere} \rhoFn(V_j(r\omega))^2\,\dd \sigma_d(\omega).
\]
Since $\rhoFn$ is $1$-Lipschitz and $U = V_j + R_j$,
\[
\rhoFn(V_j) \le \rhoFn(U) + \abs{R_j}.
\]
Squaring and using $(a+b)^2 \le 2a^2 + 2b^2$ yields
\[
\int_{\Sphere} \abs{V_j(r\omega)}^2\,\dd \sigma_d(\omega)
\le
2B^2 \int_{\Sphere} \rhoFn(U(r\omega))^2\,\dd \sigma_d(\omega)
+
2B^2 \int_{\Sphere} \abs{R_j(r\omega)}^2\,\dd \sigma_d(\omega).
\]
Now use $\abs{U}^2 \le 2\abs{V_j}^2 + 2\abs{R_j}^2$ to obtain
\[
\int_{\Sphere} \abs{U(r\omega)}^2\,\dd \sigma_d(\omega)
\le
4B^2 \int_{\Sphere} \rhoFn(U(r\omega))^2\,\dd \sigma_d(\omega)
+
(4B^2+2)\int_{\Sphere} \abs{R_j(r\omega)}^2\,\dd \sigma_d(\omega).
\]
Multiply by $\frac{2}{(d-1)!}r^{2d-1}\ee^{-r^2}$, integrate
$r \in [j,j+1]$, and sum over $j \ge 0$.  Using \eqref{eq:polar} and
Corollary~\ref{cor:leakage}, we get
\[
\norm{U}_{\Fd}^2
\le
4B^2 \norm{\rhoFn \circ U}_{\Fd}^2
+
(4B^2+2)\eta_M \norm{U}_{\Fd}^2.
\]
Absorbing the remainder term gives
\[
\norm{U}_{\Fd}^2
\le
\frac{4B_{d,M}^2}{1-(4B_{d,M}^2+2)\eta_M}
\norm{\rhoFn \circ U}_{\Fd}^2.
\]
Set
\[
C_d^2
:=
\frac{4B_{d,M}^2}{1-(4B_{d,M}^2+2)\eta_M}.
\]
Then the displayed estimate shows that
\[
\norm{U}_{\Fd}^2
\le
C_d^2\,\norm{\rhoFn \circ U}_{\Fd}^2
\]
for every holomorphic polynomial $U$ with $U(0)=0$.

Now let $U \in \Fd$ satisfy $U(0)=0$.  Let $U^{(\le N)}$ be the truncation
of its Taylor expansion to total degree at most $N$.  By
Proposition~\ref{prop:fock-facts}, these truncations satisfy
$U^{(\le N)} \to U$ in $\Fd$.  Because $U(0)=0$, the constant coefficient
of $U$ vanishes, so each $U^{(\le N)}$ also satisfies
$U^{(\le N)}(0)=0$.  Since $\rhoFn$ is
$1$-Lipschitz,
\[
\norm{\rhoFn \circ U^{(\le N)} - \rhoFn \circ U}_{\Fd}
\le
\norm{U^{(\le N)}-U}_{\Fd}
\to
0.
\]
Passing to the limit in the polynomial estimate yields
\[
\norm{U}_{\Fd}^2
\le
C_d^2\,\norm{\rhoFn \circ U}_{\Fd}^2.
\]
This proves Theorem~\ref{thm:orthog}.
\end{proof}

\begin{remark}
For the orthogonal coercivity estimate, only two genuinely new ingredients
are needed beyond the one-dimensional proof.
\begin{enumerate}[(1)]
\item Replace monomials $z^n$ by homogeneous holomorphic pieces of total
degree $n$.
\item Lift the circle estimates to the sphere by averaging the
$U(1)$-action $\omega \mapsto \ee^{\ii t}\omega$.
\end{enumerate}
Once these are in place, the block decomposition, leakage estimate, and
final absorption follow the same bookkeeping pattern as in
\texttt{Constant/Latex/FockSPR/BlockDecomposition.tex} and
\texttt{Constant/Latex/FockSPR/MainTheorem.tex}.
\end{remark}

\begin{remark}
This note does \emph{not} optimize constants.  The proof above gives a
dimension-dependent constant $C_d$.  It seems plausible that a more careful
choice of blocks and a sharper use of the one-dimensional high-frequency
estimate could keep the constants much closer to the $d=1$ case.
\end{remark}

\end{document}
